% !TEX root = ../main.tex
%% Corrigé du sujet type bac n° 2 — valeurs vérifiées : verifications/sujets_bac.py

\corrigeSujet{2}

\begin{corrigeBac}{9}{Le dihydrogène vert}
\partieExo{1}{Produire le dihydrogène}
\begin{questions}
\item Oxydation : $\ce{2H2O -> O2 + 4H+ + 4e-}$ ; réduction : $\ce{2H+ + 2e- -> H2}$, à multiplier par 2. En additionnant :
$\ce{2H2O -> 2H2 + O2}$. Le générateur impose le sens de la transformation, inverse du sens d'évolution spontanée :
la transformation est \res{\text{forcée}}. \pts{1 pt}
\item Le dihydrogène se forme par réduction, donc à la \res{\text{cathode}}, reliée à la borne $-$ du générateur : c'est
l'\res{\text{électrode B}}. Les électrons circulent dans les fils de l'électrode A vers la borne $+$, et de la borne $-$ vers
l'électrode B. \pts{0,5 pt}
\item $n(\mathrm{e}^-)=\dfrac{I\,\Delta t}{F}=\dfrac{200\times3600}{96500}=\res{\qty{7,46}{mol}}$ ; d'après la demi-équation,
$n(\ce{H2})=\dfrac{n(\mathrm{e}^-)}{2}=\qty{3,73}{mol}$, soit $m=3{,}73\times2{,}0=\res{\qty{7,46}{g}}$. \pts{1 pt}
\item $V=\dfrac{n\,R\,T}{P}=\dfrac{3{,}73\times8{,}314\times293}{\num{1.0e5}}=\qty{9.1e-2}{m^3}$, soit $\res{\qty{91}{L}}$. \pts{0,75 pt}
\item $\mathcal{E}=U\,I\,\Delta t=1{,}9\times200\times3600=\res{\qty{1,37e6}{J}}=\qty{0,38}{kW.h}$. Pour \qty{1}{kg} :
$\dfrac{0{,}38}{\num{7.46e-3}}=\res{\qty{50,9}{kW.h}}$, soit environ \qty{51}{kW.h}. \pts{1 pt}
\end{questions}
\partieExo{2}{Stocker le dihydrogène}
\begin{questions}[resume]
\item $n=\dfrac{P\,V}{R\,T}=\dfrac{\num{200e5}\times\num{50e-3}}{8{,}314\times293}=\qty{411}{mol}$, soit $m=411\times2{,}0=\res{\qty{0,82}{kg}}$. \pts{0,75 pt}
\item Le modèle néglige le volume propre des molécules et leurs interactions, ce qui n'est plus justifié lorsque le gaz
est très comprimé : les molécules sont alors proches les unes des autres. $m_\text{réelle}=\dfrac{0{,}82}{1{,}13}=\res{\qty{0,73}{kg}}$. \pts{0,75 pt}
\end{questions}
\partieExo{3}{Restituer l'énergie}
\begin{questions}[resume]
\item Anode (oxydation) : $\ce{H2 -> 2H+ + 2e-}$ ; cathode (réduction) : $\ce{O2 + 4H+ + 4e- -> 2H2O}$. Réaction de
fonctionnement : $\res{\ce{2H2 + O2 -> 2H2O}}$. \pts{0,75 pt}
\item $n(\ce{H2})=\dfrac{730}{2{,}0}=\qty{365}{mol}$ et $n(\mathrm{e}^-)=2\times365=\qty{730}{mol}$, d'où $Q=730\times96500=\res{\qty{7,04e7}{C}}$.
$\mathcal{E}=U\,Q=0{,}70\times\num{7.04e7}=\qty{4,93e7}{J}$, soit $\res{\qty{13,7}{kW.h}}$. \pts{1 pt}
\item Produire \qty{0,73}{kg} de dihydrogène a demandé $0{,}73\times50{,}9=\qty{37,2}{kW.h}$ ; la pile n'en restitue que \qty{13,7}{kW.h} :
rendement $\eta=\dfrac{13{,}7}{37{,}2}=\res{37~\%}$. Près des deux tiers de l'énergie sont perdus, sans compter l'énergie de
compression. Le dihydrogène vert reste intéressant pour stocker durablement de grandes quantités d'énergie
renouvelable, ou pour les usages où les batteries sont trop lourdes. \pts{1,5 pt}
\end{questions}
\end{corrigeBac}

\begin{corrigeBac}{6}{L'échographie Doppler}
\begin{questions}
\item Ce sont des ondes \res{\text{mécaniques}} : elles ont besoin d'un milieu matériel pour se propager. Leur fréquence,
supérieure à \qty{20}{kHz}, dépasse la limite de perception de l'oreille humaine. \pts{0,5 pt}
\item $\lambda=\dfrac{c}{f}=\dfrac{1540}{\num{5.0e6}}=\res{\qty{0,31}{mm}}$. \pts{0,5 pt}
\item Les ultrasons parcourent $2\,d$ : $\Delta t=\dfrac{2\times0{,}040}{1540}=\res{\qty{52}{\micro s}}$. \pts{0,75 pt}
\item $\Delta f>0$ : le sang se \res{\text{rapproche}} de la sonde. $v=\dfrac{c\,\Delta f}{2\,f\cos\theta}=\dfrac{1540\times1300}{2\times\num{5.0e6}\times0{,}50}=\res{\qty{0,40}{m.s^{-1}}}$. \pts{1,25 pt}
\item À \qty{5,0}{MHz}, l'atténuation vaut $0{,}5\times5{,}0=\qty{2,5}{dB}$ par centimètre ; l'aller-retour mesure \qty{8,0}{cm} :
$\res{\qty{20}{dB}}$. $10\log\dfrac{I_\text{émise}}{I_\text{reçue}}=20$ donne $\dfrac{I_\text{émise}}{I_\text{reçue}}=\res{100}$. \pts{1 pt}
\item Perpendiculairement à l'artère, $\cos\theta=0$ : le décalage serait \res{\text{nul}}, et la vitesse non mesurable. \pts{0,75 pt}
\item À \qty{2,5}{MHz}, l'atténuation vaut \qty{1,25}{dB} par centimètre : l'écho reste détectable tant que l'aller-retour ne dépasse
pas $\dfrac{60}{1{,}25}=\qty{48}{cm}$, soit une profondeur de \res{\qty{24}{cm}}, contre $\dfrac{60}{2{,}5\times2}=\qty{12}{cm}$ seulement à \qty{5,0}{MHz} :
seule la sonde de \qty{2,5}{MHz} atteint confortablement le cœur. Mais sa longueur d'onde, \qty{0,62}{mm}, est deux fois plus
grande : les détails observables sont moins fins. On choisit donc la fréquence la plus élevée compatible avec la
profondeur à explorer. \pts{1,25 pt}
\end{questions}
\schema[corrige]{bac-s2-doppler}
\end{corrigeBac}

\begin{corrigeBac}{5}{Ötzi, l'homme des glaces}
\begin{questions}
\item $\ce{^{14}_{6}C -> ^{14}_{7}N + ^{0}_{-1}e}$ (avec un antineutrino). Le nombre de nucléons ($14=14+0$) et la charge
($6=7-1$) se conservent. \pts{0,75 pt}
\item La demi-vie est la durée au bout de laquelle la moitié des noyaux initialement présents s'est désintégrée.
$\lambda=\dfrac{\ln2}{5700}=\res{\qty{1,22e-4}{an^{-1}}}$, soit $\dfrac{\num{1.22e-4}}{\num{3.16e7}}=\res{\qty{3,85e-12}{s^{-1}}}$. \pts{0,75 pt}
\item $\dfrac{A}{A_0}=\mathrm{e}^{-\lambda t}$, d'où $\ln\dfrac{A_0}{A}=\lambda\,t$ et $t=\dfrac{1}{\lambda}\ln\dfrac{A_0}{A}$. \pts{0,75 pt}
\item $t=\dfrac{1}{\num{1.22e-4}}\ln\dfrac{0{,}23}{0{,}122}$ : Ötzi est mort il y a environ $\res{\num{5.2e3}\ \text{ans}}$, vers 3200 avant notre ère. \pts{1 pt}
\item $N=\dfrac{A_0}{\lambda}=\dfrac{0{,}23}{\num{3.85e-12}}=\res{\num{6.0e10}}$ noyaux de carbone 14 par gramme de carbone. \pts{0,75 pt}
\item Après \num{50000} ans, environ 8,8 demi-vies, l'activité ne vaut plus que $2^{-50000/5700}=0{,}23~\%$ de $A_0$, soit
\qty{5,3e-4}{Bq} par gramme : elle devient trop faible pour être distinguée de la radioactivité ambiante. \pts{1 pt}
\end{questions}
\end{corrigeBac}
